\begin{table*}[t]
\caption{Conditional standard-fit comparison. $N_{\rm BH}^{\rm ref}$ is the expected count in the selected maximum-quadrature-weight joint row. Published intervals are empirical 16th--84th percentiles of 4,000 released derived model realizations per cluster. $f_D$ is the BH contribution to the local diffusion trace at zero remnant velocity, with a common constant Coulomb logarithm. The probe mass is $60\,M_\odot$. These are model quantities, not directly observed counts or residence fractions. The NGC 6397 standard fit is a sensitivity case; its preferred no-BH chain is unavailable.}\label{tab:cluster-transport-comparison}
\centering\begin{tabular}{lrrrrr}
\toprule
Cluster & $N_{\rm BH}^{\rm ref}$ & Published $N_{\rm BH}$ & $m_{\rm r}/M_{\rm BH}^{\rm ref}$ & $f_D(0)$ (\%) & $f_D(0.1r_h)$ (\%) \\
\midrule
47 Tuc & 65.9 & 61.9 [51.2, 81.1] & 0.1357 & 81.579 & 0.01328 \\
NGC 1851 & 16 & 15 [13.1, 17.6] & 0.3034 & 99.542 & 0.0004788 \\
NGC 3201 & 4.28 & 0.926 [0, 4.13] & 0.981 & 95.836 & 0.0846 \\
$\omega$ Cen & 12,100 & 12,400 [12,100, 13,000] & 0.0003391 & 99.999 & 98.52 \\
M4 & 5.33 & 6.17 [3.99, 9.65] & 0.8782 & 99.431 & 0.02131 \\
NGC 6397 (standard) & 2.94 & 3.08 [2.84, 3.31] & 3.041 & 97.793 & $<10^{-4}$ \\
\bottomrule
\end{tabular}
\end{table*}
